\documentclass[11pt]{article}
\usepackage[margin=25mm]{geometry}
\usepackage{amsmath,amssymb}
\begin{document}

\title{I01 prevalidation 12 protocol:\\Exact atlas of paired-reference leakage conditions}
\date{2026-10-04}
\maketitle

\section*{Question}
Stages 9--11 show that one paired $R_x$ intervention can make a previously exact
$P=0$ pooled null observable. This protocol asks which circuit features are
necessary: two-sided placement, sector-mixing axis, same-qubit support, or inverse
sign.

\section*{Frozen exact-Aer conditions}
At $\epsilon=1.0$ rad and workers 6, all two blocks, three axes, seven phase labels,
both $\tau\in\{0,\pi/2\}$, and unrestricted/$P=0$ conditions are evaluated for six
rows: one-sided $R_x$ on qubit zero; same-qubit paired inverse $R_x$; same-qubit
paired inverse $R_y$; same-qubit paired inverse $R_z$; inverse $R_x$ from qubit zero
to qubit one; and same-qubit equal-sign $R_x$. The first operation follows
preparation; the second follows the science layer before reference projection.

\section*{Decision discipline}
This stage makes no finite-shot significance claim. It records full pooled contrast
tables, especially the $P=0$, $\tau=0$ constraint row. A later sampling experiment
may select a condition only after this complete map is observed and its sampling
criterion is written down.

\section*{Reproduction}
\texttt{../src/relational\_atlas\_i01/holdout/qiskit\_rank\_holdout12\_intervention\_atlas.py}
\end{document}
